\section{Prapemrosesan Data \& Parameter Statistik}

\begin{frame}
\centering
\highlight{\Huge Bagian 2: Prapemrosesan Data \& Parameter Statistik}
\end{frame}

\subsection{Akuisisi \& Log-Return}
\begin{frame}{2.1 Akuisisi \& Prapemrosesan Data Mentah}
Contoh Dataset (1-3 Nov 2023): ADRO, BBCA, ICBP, TLKM
\vspace{0.2cm}

\textbf{Harga Penutupan (\textit{Close Prices}):}
\begin{table}[]
\resizebox{\textwidth}{!}{
\begin{tabular}{lrrrr}
\toprule
Date & ADRO.JK & BBCA.JK & ICBP.JK & TLKM.JK \\
\midrule
2023-11-01 & 2681.33 & 8783.50 & 10173.34 & 3465.73 \\
2023-11-02 & 2658.98 & 9035.21 & 10148.81 & 3465.73 \\
\bottomrule
\end{tabular}
}
\end{table}

\textbf{Transformasi Imbal Hasil (\textit{Log-Return}):}
\begin{equation}
R_{i,t} = \ln \left( \frac{P_{i,t}}{P_{i,t-1}} \right)
\end{equation}

\begin{table}[]
\begin{tabular}{lrrrr}
\toprule
Date & ADRO & BBCA & ICBP & TLKM \\
\midrule
2023-11-02 & -0.0083 & 0.0286 & -0.0024 & 0.0000 \\
2023-11-03 & 0.0409 & 0.0056 & 0.0096 & -0.0139 \\
\bottomrule
\end{tabular}
\end{table}
\end{frame}

\subsection{Estimasi Parameter Statistik}
\begin{frame}{2.2 Estimasi Parameter Statistik Empiris}
Hasil kalkulasi parameter statistik dasar berdasarkan data sampel (2 hari):

\begin{table}[]
\resizebox{\textwidth}{!}{
\begin{tabular}{lrrrr}
\toprule
Parameter & ADRO & BBCA & ICBP & TLKM \\
\midrule
$\bar{R}_i$ & 0,0163 & 0,0171 & 0,0036 & -0,0070 \\
$\text{std}(R_i)$ & 0,0348 & 0,0163 & 0,0085 & 0,0098 \\
$\mu_{ann,i}$ & 4,1076 & 4,3092 & 0,9072 & -1,7514 \\
$\sigma_{ann,i}$ & 0,5522 & 0,2582 & 0,1347 & 0,1560 \\
\bottomrule
\end{tabular}
}
\end{table}

\textbf{Aggregated Sharpe Ratio ($Z$):}
\begin{itemize}
    \item Mean $|\mu_{ann}| = \frac{1}{4}(4,11 + 4,31 + 0,91 + 1,75) = 2,7689$
    \item Mean $\sigma_{ann} = \frac{1}{4}(0,55 + 0,26 + 0,13 + 0,16) = 0,2753$
    \item $Z = 2,7689 / 0,2753 = \mathbf{10,0585}$
\end{itemize}
\end{frame}

\subsection{Risk-Aversion Endogen}
\begin{frame}{2.3 Formulasi Risk-Aversion Endogen ($\gamma$)}
Parameter $\gamma$ merespons kondisi pasar secara dinamis (*endogenous*) melalui fungsi sigmoid logistik.

\textbf{Rumus Risk-Aversion:}
\begin{equation}
\gamma = \frac{1}{1 + e^{Z}}
\end{equation}

\textbf{Hasil Kalkulasi Sampel:}
\begin{itemize}
    \item Dengan $Z = 10,0585$, diperoleh $\gamma \approx 4,28 \times 10^{-5}$.
    \item \textit{Note:} Pada implementasi \textit{backtesting} (63 hari), karakteristik pasar yang lebih stabil menghasilkan $Z \approx 3,43$ sehingga diperoleh $\gamma = \mathbf{0,031156}$.
\end{itemize}
\end{frame}

\subsection{Binarisasi \& Imbal Hasil Strategis}
\begin{frame}{2.4 Binarisasi \& Imbal Hasil Strategis ($\tilde{\mu}$)}
Data binarisasi digunakan untuk menangkap \textit{microstates} sistem: $s_{i,t} = 1$ jika $R_{i,t} \leq 0$, dan $s_{i,t} = 0$ jika $R_{i,t} > 0$.

\textbf{Imbal Hasil Strategis Bersyarat ($\tilde{\mu}_i$):}
\begin{equation}
\tilde{\mu}_i = \sum_{s \in S} P(s) \bar{R}_i(s)
\end{equation}

\textbf{Hasil Vektor $\tilde{\mu}$ (Final):}
\begin{equation}
\tilde{\mu} = \begin{bmatrix} 0,0041 &  0,0008 &  0,0031 & -0,0037 \end{bmatrix}^T
\end{equation}
\end{frame}
